% =============================================================================
%  TEMPLATE PubMed / BIOMÉDICO REFINADO v1.0
%  Estrutura IMRAD (Introdução, Métodos, Resultados e Discussão)
%  Estilo Vancouver (citações numéricas) — padrão PubMed/Medline
%  Compilar: pdflatex -> pdflatex
% =============================================================================
\documentclass[12pt,a4paper]{article}

% --- Pacotes ---
\usepackage[utf8]{inputenc}
\usepackage[T1]{fontenc}
\usepackage{amsmath,amssymb}
\usepackage{graphicx}
\usepackage[hmargin=2.5cm,vmargin=2.5cm]{geometry}
\usepackage[brazil]{babel}
\usepackage{xcolor}
\usepackage{hyperref}
\hypersetup{
    colorlinks=true,
    linkcolor=blue!60!black,
    citecolor=red!60!black,
    urlcolor=blue!60!black,
    breaklinks=true,
    pdfauthor={Autor},
    pdftitle={Título do Artigo Biomédico}
}
\usepackage[numbers]{natbib}  % Vancouver numérico
\bibliographystyle{unsrtnat}
\usepackage{booktabs}
\usepackage{caption}
\usepackage{enumitem}
\usepackage{array}
\usepackage{microtype}
\usepackage{times}            % Fonte Times New Roman (padrão biomédico)

% --- Comandos auxiliares ---
\newcommand{\ICMJE}{\textit{ICMJE}}
\newcommand{\IC}{\textit{intervalo de confian\c{c}a}}
\newcommand{\OR}{\textit{odds ratio}}
\newcommand{\RR}{\textit{risco relativo}}

% =============================================================================
\begin{document}

% --- Título ---
\title{T\'itulo do Artigo: Subt\'itulo Explicativo \\
        \large Template Refinado para Peri\'odicos Indexados no PubMed}

% --- Autores e afiliações ---
\author{Marcos A. Santos$^{1*}$\thanks{Correspond\^encia: marcos.santos@exemplo.com},
        Ana C. Oliveira$^{2}$,
        Carlos R. Lima$^{1}$ \\
        \small $^{1}$Departamento de Medicina, Universidade Federal do Cear\'a,
        Fortaleza, CE, Brasil \\
        \small $^{2}$Instituto de Sa\'ude Coletiva, Universidade de S\~ao Paulo,
        S\~ao Paulo, SP, Brasil \\
        \small $^{*}$Autor correspondente}

\date{\today}

\maketitle

% =============================================================================
\begin{abstract}
\textbf{Objetivo}: Demonstrar a estrutura de um artigo biom\'edico
refinado para submiss\~ao a peri\'odicos indexados no PubMed,
seguindo as recomenda\c{c}\~oes do \ICMJE{} e o padr\~ao Vancouver
de cita\c{c}\~oes.

\textbf{M\'etodos}: Este template foi elaborado em \LaTeX{} com
a classe \texttt{article}, estrutura IMRAD e refer\^encias no
formato num\'erico (Vancouver). O c\'odigo-fonte est\'a preparado
para compila\c{c}\~ao direta com \texttt{pdflatex}.

\textbf{Resultados}: O template apresenta todas as se\c{c}\~oes
necess\'arias para um artigo biom\'edico completo, incluindo tabelas
de dados cl\'inicos, an\'alise estat\'istica e refer\^encias
formatadas segundo o estilo Vancouver.

\textbf{Conclus\~ao}: Este template oferece uma base s\'olida para
submiss\~ao a peri\'odicos biom\'edicos, reduzindo o tempo de
formata\c{c}\~ao e garantindo conformidade com os padr\~oes
internacionais.

\textbf{Palavras-chave}: template LaTeX; PubMed; estilo Vancouver;
ICMJE; artigo biom\'edico.
\end{abstract}

% =============================================================================
\section{Introdu\c{c}\~ao}
\label{sec:intro}

A publica\c{c}\~ao em peri\'odicos indexados no PubMed exige
conformidade com diretrizes rigorosas de formata\c{c}\~ao e
estrutura editorial \cite{ICMJE2024}. O estilo Vancouver de
cita\c{c}\~oes, caracterizado por refer\^encias num\'ericas
ordenadas por aparecimento, \'e o padr\~ao adotado pela maioria
das revistas biom\'edicas \cite{Patrias2007Citing}.

Estima-se que aproximadamente 30 milh\~oes de artigos estejam
indexados no PubMed, com crescimento anual de cerca de 5\%
\cite{Landhuis2024PubMed}. Apesar da ampla disponibilidade de
templates institucionais, muitos pesquisadores ainda enfrentam
dificuldades com a formata\c{c}\~ao manual de manuscritos.

O objetivo deste trabalho \'e fornecer um template \LaTeX{}
refinado, completo e de f\'acil adapta\c{c}\~ao, seguindo as
recomenda\c{c}\~oes do \ICMJE{} e do Comit\^e Internacional de
Editores de Peri\'odicos M\'edicos.

% =============================================================================
\section{M\'etodos}
\label{sec:metodos}

\subsection{Desenho do Estudo}

Trata-se de um template metodol\'ogico, do tipo guia de
formata\c{c}\~ao, elaborado conforme as diretrizes do \ICMJE{}
\cite{ICMJE2024} e adaptado para a realidade da pesquisa
biom\'edica brasileira.

\subsection{Estrutura do Template}

O template segue a estrutura IMRAD:

\begin{enumerate}[label=\textbf{\arabic*.}]
    \item \textbf{Introdu\c{c}\~ao}: contextualiza\c{c}\~ao e
        objetivos do estudo;
    \item \textbf{M\'etodos}: descri\c{c}\~ao detalhada dos
        procedimentos;
    \item \textbf{Resultados}: apresenta\c{c}\~ao objetiva dos
        dados;
    \item \textbf{Discuss\~ao}: interpreta\c{c}\~ao e compara\c{c}\~ao
        com a literatura.
\end{enumerate}

\subsection{An\'alise Estat\'istica}

Os dados hipot\'eticos foram analisados utilizando o software R
vers\~ao 4.3 (R Foundation for Statistical Computing). As vari\'aveis
cont\'inuas foram expressas como m\'edia $\pm$ desvio padr\~ao
(DP) ou mediana e amplitude interquartil (AIQ), conforme a
distribui\c{c}\~ao. Vari\'aveis categ\'oricas foram expressas como
frequ\^encia absoluta ($n$) e relativa (\%). O teste $t$ de Student
foi utilizado para compara\c{c}\~ao de m\'edias e o teste
$\chi^2$ para propor\c{c}\~oes. Valores de $p < 0.05$ foram
considerados estatisticamente significativos.

% =============================================================================
\section{Resultados}
\label{sec:resultados}

\subsection{Caracter\'isticas da Amostra}

A amostra hipot\'etica foi composta por 250 pacientes, sendo
140 (56\%) do sexo feminino e 110 (44\%) do sexo masculino, com
idade m\'edia de 52.3 $\pm$ 14.7 anos (Tabela~\ref{tab:demografica}).

\begin{table}[htbp]
\centering
\caption{Caracter\'isticas demogr\'aficas e cl\'inicas da
         amostra $(n = 250)$}
\label{tab:demografica}
\begin{tabular}{lcc}
\toprule
Vari\'avel & Grupo Controle $(n = 125)$ & Grupo Estudo $(n = 125)$ \\
\midrule
Idade (m\'edia $\pm$ DP) & 51.8 $\pm$ 15.2 & 52.7 $\pm$ 14.1 \\
Sexo feminino, $n$ (\%) & 68 (54.4\%) & 72 (57.6\%) \\
IMC (kg/m$^2$) & 26.8 $\pm$ 4.3 & 27.2 $\pm$ 4.6 \\
Press\~ao arterial sist\'olica (mmHg) & 128.4 $\pm$ 16.7 & 132.1 $\pm$ 18.2$^*$ \\
Glicemia de jejum (mg/dL) & 94.6 $\pm$ 12.3 & 98.1 $\pm$ 15.7$^*$ \\
\midrule
\multicolumn{3}{l}{\small $^* p < 0.05$ vs. controle} \\
\end{tabular}
\end{table}

\subsection{Desfecho Prim\'ario}

O desfecho prim\'ario ocorreu em 38 pacientes do grupo estudo
(30.4\%) versus 52 do grupo controle (41.6\%), com
$\RR = 0.73$ (\IC{} 95\%: 0.52--0.98; $p = 0.038$),
conforme ilustrado na Figura~\ref{fig:kaplan}.

\begin{figure}[htbp]
\centering
\vspace{2cm}
\begin{minipage}{0.75\textwidth}
\centering
\fbox{\parbox{\textwidth}{\centering
    [Curva de Kaplan-Meier: sobrevida livre de evento]}}
\end{minipage}
\caption{Curva de Kaplan-Meier para o desfecho prim\'ario.}
\label{fig:kaplan}
\end{figure}

\subsection{An\'alise Multivariada}

A an\'alise de regress\~ao log\'istica multivariada identificou
os seguintes preditores independentes (Tabela~\ref{tab:regressao}):

\begin{table}[htbp]
\centering
\caption{Regress\~ao log\'istica multivariada para o desfecho
         prim\'ario}
\label{tab:regressao}
\begin{tabular}{lccc}
\toprule
Vari\'avel & $\OR$ & IC 95\% & $p$ \\
\midrule
Interven\c{c}\~ao (grupo estudo) & 0.68 & 0.48--0.95 & 0.024 \\
Idade (por 10 anos) & 1.24 & 1.08--1.42 & 0.003 \\
IMC (por 5 kg/m$^2$) & 1.31 & 1.12--1.53 & 0.001 \\
Sexo masculino & 1.15 & 0.82--1.61 & 0.418 \\
\bottomrule
\end{tabular}
\end{table}

% =============================================================================
\section{Discuss\~ao}
\label{sec:discussao}

Este template demonstra a estrutura completa de um artigo
biom\'edico no formato exigido por peri\'odicos indexados no
PubMed. Os resultados hipot\'eticos apresentados ilustram a
aplica\c{c}\~ao dos principais elementos estat\'isticos e
gr\'aficos utilizados na literatura m\'edica.

A redu\c{c}\~ao de 27\% no risco relativo observada no grupo
estudo \'e consistente com achados de ensaios cl\'inicos
pr\'evios \cite{Topol2024Preventive}. A an\'alise multivariada
confirmou a independ\^encia do efeito ap\'os ajuste para
vari\'aveis confundidoras.

Limita\c{c}\~oes deste template incluem o car\'ater ilustrativo
dos dados e a aus\^encia de valida\c{c}\~ao externa.
Pesquisadores devem adaptar o conte\'udo e as an\'alises
\`a sua \'area espec\'ifica de estudo.

% =============================================================================
\section{Conclus\~ao}
\label{sec:conclusao}

Apresentamos um template \LaTeX{} refinado e completo para
submiss\~ao a peri\'odicos indexados no PubMed. A estrutura
IMRAD, as cita\c{c}\~oes no estilo Vancouver e a formata\c{c}\~ao
segundo as recomenda\c{c}\~oes do \ICMJE{} garantem conformidade
com os padr\~oes editoriais da \'area biom\'edica.

% =============================================================================
% Agradecimentos
\section*{Agradecimentos}
Os autores agradecem ao suporte financeiro da CAPES
(C\'odigo de Financiamento 001) e do CNPq (Processo
123456/2024-0).

% =============================================================================
% Conflito de interesses
\section*{Conflito de Interesses}
Os autores declaram n\~ao haver conflito de interesses.

% =============================================================================
% Contribuição dos autores
\section*{Contribui\c{c}\~ao dos Autores}
\textbf{MAS}: concep\c{c}\~ao do estudo, reda\c{c}\~ao do manuscrito;
\textbf{ACO}: an\'alise estat\'istica, revis\~ao cr\'itica;
\textbf{CRL}: orienta\c{c}\~ao, revis\~ao final.

% =============================================================================
% Referências no estilo Vancouver
\begin{thebibliography}{99}

\bibitem{ICMJE2024}
International Committee of Medical Journal Editors.
\textit{Recommendations for the Conduct, Reporting, Editing,
and Publication of Scholarly Work in Medical Journals}.
ICMJE, 2024. Dispon\'ivel em: \url{https://www.icmje.org/}

\bibitem{Patrias2007Citing}
K. Patrias.
\textit{Citing Medicine: The NLM Style Guide for Authors,
Editors, and Publishers}. 2nd ed. Bethesda (MD):
National Library of Medicine, 2007.

\bibitem{Landhuis2024PubMed}
E. Landhuis.
PubMed at 30: a look back at the biomedical database that
changed science. \textit{Nature}. 2024;626(7998):446--448.

\bibitem{Topol2024Preventive}
E.J. Topol.
Preventive cardiology in the era of precision medicine.
\textit{Lancet}. 2024;403(10428):912--925.

\end{thebibliography}

\end{document}
